\documentclass[11pt]{article}
\usepackage[margin=0.8in]{geometry}
\usepackage{enumitem}
\usepackage{setspace}
\setlist[itemize]{leftmargin=1.15em, itemsep=0.1em, topsep=0.18em}
\setlist[enumerate]{leftmargin=1.35em, itemsep=0.1em, topsep=0.18em}
\setlength{\parindent}{0pt}
\setlength{\parskip}{0.26em}
\setstretch{1.02}
\pagestyle{plain}

\begin{document}

\begin{center}
{\Large\bfseries Uncertainty and Information Economics}\\[0.3em]
Sample syllabus\\[0.15em]
Ke Cheng HSU\\
Department of Economics, National Chengchi University\\
Undergraduate elective, 3 credits, one semester, taught in English
\end{center}

\textbf{Place in the curriculum.}
This is the department's undergraduate elective on uncertainty and information. The first part is decision under risk and the market allocation of risk. The second part is strategic interaction when people know different things: auctions, information in markets, mechanism design, and moral hazard. Industrial Economics treats price discrimination as a pricing tool. The mechanism-design weeks here ask what a contract can implement when the buyer knows something the seller does not.

\textbf{Prerequisites.}
Intermediate microeconomics. Basic calculus and algebra. A prior game-theory course is not required. The equilibrium concepts are taught in the weeks that use them.

\textbf{Readings.}
Lecture notes are the text. They are posted before each week. The following books cover the same material and are optional.

\begin{itemize}
\item Danthine, Jean-Pierre, and John B.\ Donaldson. \emph{Intermediate Financial Theory}. Academic Press, 2014. The weeks on contingent claims and asset pricing.
\item Tadelis, Steven. \emph{Game Theory: An Introduction}. Princeton University Press, 2013. The weeks on Bayesian games and auctions.
\item Macho-Stadler, In\'es, and J.\ David P\'erez-Castrillo. \emph{An Introduction to the Economics of Information: Incentives and Contracts}. The weeks on moral hazard and on second-best mechanisms.
\item Bolton, Patrick, and Mathias Dewatripont. \emph{Contract Theory}. MIT Press, 2005. A further reference for a student who continues to graduate work.
\end{itemize}

\textbf{How the course runs.}
Three hours a week. Two hours are lecture. The third hour takes up the harder questions on the problem set. Six problem sets. The lowest score is dropped. Group discussion is welcome. Each student submits their own written answers. A serious attempt at the setup earns credit when the last step is wrong. Solutions are posted after the due date.

\textbf{Learning outcomes.}
By the end of the semester, students should be able to
\begin{itemize}
\item use expected utility, risk aversion, and the certainty equivalent,
\item describe how a market for contingent claims shares risk, and state the logic of a consumption-based asset price,
\item find a Bayesian Nash equilibrium in a simple auction or bargaining problem,
\item explain the winner's curse, and why a herd can form from rational updating,
\item state the revelation principle and write the risk-incentive tradeoff in a moral-hazard contract.
\end{itemize}

\textbf{Assessment.}
Problem sets 20 percent. Midterm examination 30 percent, in class, covering weeks 1--5. A documented absence on the midterm moves that weight to the final. Final examination 50 percent, cumulative, with more weight on the information weeks. The questions have the same form as the problem sets. Closed book.

\textbf{Weekly plan.}
\begin{enumerate}
\item Expected utility. Lotteries, the independence axiom in one example, and expected value against expected utility.
\item Risk attitude and optimal risk bearing. Concavity, the certainty equivalent, and the risk premium.
\item A pure-exchange economy with contingent claims. Full insurance at fair odds.
\item The asset market. A risky asset and a risk-free asset. Who holds the risk.
\item Consumption-based asset pricing. The price of an asset as expected payoff weighted by marginal utility. A mean-variance remark in the spirit of the CAPM.
\item Midterm examination on weeks 1--5.
\item Bayesian Nash equilibrium. Types, beliefs, and best responses.
\item Auctions and bargaining under incomplete information. A second-price auction with private values. One bargaining split.
\item Belief updating. Bayes' rule in a numerical signal problem.
\item Common-value auctions and the winner's curse. Voting as a second application of the same updating, if time allows.
\item Information cascades. A rational herd. Perfect Bayesian equilibrium in one reputation game.
\item The revelation principle. A direct mechanism and why it is enough to study truthful ones.
\item An efficient mechanism. The pivot mechanism in a public-good example.
\item A second-best mechanism. Incentive compatibility and a distorted contract for the low type.
\item Moral hazard. Hidden effort, a participation constraint, and an incentive constraint with two effort levels.
\item The risk-incentive tradeoff. Why the agent is not fully insured. A remark on multitasking and on career concern.
\item Review.
\item Final examination.
\end{enumerate}

\end{document}
